% RQ3 Abstract: Proposal Pipeline Effects

The proposal pipeline---the sequence of stages a governance proposal traverses from inception to execution---fundamentally shapes DAO decision-making. Pipeline design involves tradeoffs between throughput (processing proposals quickly) and quality (ensuring adequate deliberation and filtering).

We investigate how proposal pipeline settings affect governance outcomes using multi-agent simulation. Through \experimentruns{} simulation runs, we examine three key pipeline mechanisms: (1) temp-checks that filter low-support proposals before formal voting, (2) fast-track provisions that expedite uncontroversial or urgent proposals, and (3) expiry windows that bound deliberation time.

Our results reveal significant throughput-quality tradeoffs. Aggressive temp-check thresholds reduce formal voting load but risk filtering viable proposals with initially low awareness. Fast-tracks improve time-to-decision for qualifying proposals but may bypass deliberation on consequential issues. Tight expiry windows prevent stale proposals but can cause abandonment of complex initiatives.

We provide parameter guidelines for pipeline design and identify configurations that balance efficiency with governance quality. Our findings inform practitioners designing proposal workflows for operational DAOs.
